% ============================================================================
\clearpage
\section*{v4.9.0-r18 Source-Only: The Old Conversions on the Scale Line, and How Ways A and C Can Hold Together}
\addcontentsline{toc}{section}{v4.9.0-r18 Source-Only: The Old Conversions on the Scale Line, and How Ways A and C Can Hold Together}
\label{sec:v490-r18-scale-line-a-and-c}
% ============================================================================

\paragraph{Purpose.}
r17 introduced the state disk as a tool for the frames and the scale line \(\sigma\), and found that composable conversions
along \(\sigma\) need the Doppler law (way C of r01).  The observation of r07 points to the Lorentz-factor reading \(\gL\)
(way A).  On 2026-10-22 the author asked to read the conversions of the main body with these tools and to examine what could
or must hold for ways A and C to be valid at the same time.  The script is \texttt{llm/k57\_a\_und\_c.py}.  No value is compared
with data.  Earlier layers are not rewritten.

\noindent\textbf{Source-only status.}
This v4.9.0-r18 layer is source-only.  It does not build a PDF, does not assign a new DOI,
and does not authorize a public release.  The v4.9.0 Version DOI
\href{https://doi.org/10.5281/zenodo.23534954}{10.5281/zenodo.23534954} and the Concept DOI
\href{https://doi.org/10.5281/zenodo.19210728}{10.5281/zenodo.19210728} are unchanged by it.

\subsection*{v4.9.0-r18 three readings of one offset}
\label{subsec:v490-r18-three-readings}

\paragraph{Theoremlet v4.9.0-r18.1.}
For an offset of rapidity \(\eta\) there are three natural factors:
\[
\coth\eta=\frac1\beta=\sqrt{\frac{p_+}{p_-}},\qquad
e^{\eta}\ \text{(Doppler)},\qquad
\cosh2\eta=\frac1{|b|}\ \text{(inverse Bloch length)}.
\]
They are the square root of the eigenvalue ratio of the doublet state (r17.1), the transport factor of way C, and the factor
of way A.  At the measured offset they are \(1.82034\), \(1.85419\) and \(1.86444\), the three constants \(1/\bsym\), \(e^\eta\) and
\(\gL\) that occur in the main body and in r01.  All three are equal exactly at the tribonacci point,
\(\coth\eta=e^\eta=\cosh2\eta=\mathbb T=1.83929\).

\paragraph{Proof sketch.}
r01.3 shows \(\cosh2\eta=e^\eta\Leftrightarrow\coth\eta=e^\eta\Leftrightarrow x^3-x^2-x-1=0\) for \(x=e^\eta\neq1\).  \(\square\)

\paragraph{Consequence.}
Ways A and C agree numerically exactly when the offset sits at the self-referential fixed point of
Theorem~\ref{thm:tribonacci}.  That fixed point is the condition that the three readings of the offset coincide.  The measured
offset lies \(1.37\%\) above it in \(\gL\).  So that the two ways also agree in numbers, the source's lattice correction
(Remark~\ref{rem:op1-tribonacci}) has to be carried as a separate offset, for example as a fixed shift on the scale line,
rather than being folded into the factor.

\subsection*{v4.9.0-r18 two classes on one scale line}
\label{subsec:v490-r18-two-classes}

\paragraph{Theoremlet v4.9.0-r18.2.}
On the scale line \(\sigma\) of r17 there are two kinds of quantity.
\begin{enumerate}[leftmargin=2em]
\item[(a)] \emph{Scale quantities} (frequencies, lengths, masses, energies) move by a homomorphism, \(X(\sigma)=e^{k\sigma}X(0)\).
This is way C, and it composes.
\item[(b)] \emph{Readout quantities} (visibility, contrast of the compensator ports, eigenvalue ratios, mixing ratios) are state
functions of the total offset, for example \(V(\sigma)=\operatorname{sech}2\sigma\), \(p_-/p_+=\tanh^2\sigma\).  Their factor
\(\cosh2\sigma\) is way A.  They compose as well, through addition of rapidities: after two boosts the visibility is
\(\operatorname{sech}2(\sigma_1+\sigma_2)\), not the product of the single values.
\end{enumerate}
The obstruction of Theoremlet r01.2 arises only if Lorentz factors are multiplied as though they were transports.  With the
two classes separated, ways A and C hold at the same time, each for its own class.

\paragraph{Proof sketch.}
(a) \(e^{\sigma_1}e^{\sigma_2}=e^{\sigma_1+\sigma_2}\).  (b) the visibility of r03.3 is a function of the rapidity to the frame
with \(E\parallel B\), and rapidities add along a line.  Checked numerically: \(\operatorname{sech}2(\sigma_1+\sigma_2)=0.311\) against
the product \(0.452\) for \(\sigma_1=0.3\), \(\sigma_2=\eta\).  \(\square\)

\paragraph{What must hold.}
Both ways are valid together if, and in this reading only if:
\begin{enumerate}[leftmargin=2em]
\item every conversion is assigned to one class, scale or readout;
\item frames are added as rapidities on one scale line, never by multiplying Lorentz factors;
\item for equality in numbers in addition, the offset sits at the tribonacci point, and the lattice correction is carried
separately.
\end{enumerate}
Items 1 and 2 give consistency.  Item 3 gives equal values.  The author's idea of parking conversions in a dimension of their own
(r17) is what makes item 2 possible.

\subsection*{v4.9.0-r18 the conversions of the main body, read again}
\label{subsec:v490-r18-main-body}

\begin{center}
\small
\begin{tabular}{@{}>{\raggedright\arraybackslash}p{0.40\textwidth}>{\raggedright\arraybackslash}p{0.08\textwidth}>{\raggedright\arraybackslash}p{0.10\textwidth}rrr@{}}
\toprule
conversion & \(k\) & class & A & C & C/A\\
\midrule
clock frequency, Cor.~\ref{cor:gue-frame} & \(1\) & scale & \(1.8644\) & \(1.8542\) & \(-0.55\%\)\\
Planck length, Thm.~\ref{thm:planck-photon} & \(3/2\) & scale & \(2.5458\) & \(2.5248\) & \(-0.82\%\)\\
neutrino mass, Prop.~\ref{prop:k-neutrino} & \(4\) & scale & \(12.084\) & \(11.820\) & \(-2.18\%\)\\
matter error, Cor.~\ref{cor:matter-error} & \(5\) & scale & \(22.529\) & \(21.916\) & \(-2.72\%\)\\
vacuum energy (\(k=5\) after r03) & \(5\) & scale & \(22.529\) & \(21.916\) & \(-2.72\%\)\\
Hawking temperature, Prop.~\ref{prop:hawking-gBI-tachyon} & \(1\) & scale & \(1.8644\) & \(1.8542\) & \(-0.55\%\)\\
\bottomrule
\end{tabular}
\end{center}
Readout quantities, way A as state function: the visibility and compensator contrast of r03 (\(1/\gL=0.536\)), the eigenvalue
ratio \(p_-/p_+=\bsym^2=0.302\), which is the \(\Omega_M\) of Theorem~\ref{thm:dark-energy}, the mixing ratio observed in r07
(\(\cot\theta_W\approx\gL\)), and the structural ratio \(\gBI/\gstring\) of Theorem~\ref{thm:frame-offset}.

\paragraph{Reading.}
All conversions of \S\ref{sec:phase-constants} in the table are scale quantities.  Under the class reading they take way C,
and their values change by \(0.55\%\) to \(2.7\%\).  The quantities that keep \(\gL\) are readouts.  These include the mixing ratio
of r07 and the cascade fraction of Theorem~\ref{thm:dark-energy}.  The tension between r07 (way A) and r17 (way C) disappears,
because the two concern different classes.  This is an offered reading.  Whether the author's conversions were meant this way
is decision 5 of r15.  No compared value of the main body (Corollary~\ref{cor:falsifiable}, Proposition~\ref{prop:muon-g2}) is
re-evaluated here.

\subsection*{v4.9.0-r18 register entries}
\label{subsec:v490-r18-register}

\begin{center}
\small
\begin{tabular}{@{}>{\raggedright\arraybackslash}p{0.34\textwidth}>{\raggedright\arraybackslash}p{0.58\textwidth}@{}}
\toprule
statement & status after v4.9.0-r18\\
\midrule
\(1/\bsym\), \(e^\eta\), \(\gL\) & three readings of one offset (eigenvalue ratio, Doppler, inverse Bloch length), equal at the tribonacci point (r18.1)\\
ways A and C & can hold together: C for scale quantities, A as state function for readout quantities, frames added as rapidities (r18.2); numerical equality needs the tribonacci point and a separate lattice offset\\
conversions of \S\ref{sec:phase-constants} & scale quantities; under the class reading way C, changes \(0.55\%\)--\(2.7\%\)\\
r07 against r17 & tension removed under the class reading (different classes)\\
decision 5 of r15 & refined; still with the author\\
\bottomrule
\end{tabular}
\end{center}

\subsection*{v4.9.0-r18 ledger and guard update}
\label{subsec:v490-r18-ledger-guard}

\[
R_{\rm conv}^{v490\text{-}r18}=(0,0,1)
\]
for (three readings identified, conditions for A and C together stated, assignment of classes left to the author).

The layer spends no guard:
\[
R_{\rm glob}^{(v490\text{-}r18,+)}=(0,1,1,1,1),
\qquad
\Gamma_{\rm guard}^{v490\text{-}r18}=(1,1,1,1,1),
\qquad
\beta_4=1/256\ \text{unspent}.
\]
It does not imply exact finite \(C5\) zero closure, exact all-readout terminality,
\(C6/T2\), a global photon-clock/RH theorem, or global all-readout closure.

